
Batch Normalization exhibits 9.41 percentage points higher worst-group accuracy gap than Layer Normalization at matched average accuracy on proof-of-concept synthetic spurious correlation data (p < 0.001, d = 3.94). This effect operates via gradient flow asymmetry: Batch Normalization shows 26\% higher gradient magnitude toward spurious-aligned samples during early training.

These findings are based on synthetic data. Real dataset validation (Waterbirds, CelebA) is required before claiming generalization to real-world spurious correlation tasks. The planned attention mechanism hypothesis (CBAM, ViT) was incomplete due to technical failures and is deferred to future work.

Future work includes:
\begin{enumerate}
\item Real dataset validation on Waterbirds and CelebA
\item Completion of attention mechanism hypothesis (h-m2)
\item Optimizer ablation (SGD, Adam, AdamW)
\item Learning rate schedule ablation
\item Cross-domain evaluation (NLP, audio, tabular)
\end{enumerate}

Accuracy-matched temporal comparison provides a method for isolating architectural effects during training. Layer Normalization's instance-level normalization eliminates batch-level spurious signal aggregation present in Batch Normalization.
